%% -*- coding:utf-8 -*-
\chapter{Phrasenstrukturgrammatiken und \xbart}
\label{chap-psg-xbar}\label{chap-psg}
\lehead{\thechapter~~Phrasenstrukturgrammatiken und $\skoverline{X}$ Theory}

Dieses Kapitel führt Phrasenstrukturgrammatiken (PSGs) ein, die in vielen seit \citew{Chomsky57a}
entwickelten Theorien eine wichtige Rolle spielen. Die in diesem Kapitel entwickelte
Phrasenstrukturgrammatik wird die Grundlage für kompliziertere Phänomene bilden, die in den folgenden
Kapiteln behandelt werden. Dieses Kapitel beschäftigt sich hauptsächlich mit dem \ili{Deutsch}en und dem
\ili{Englisch}en, was für die Einführung des formalen Apparats von Phrasenstrukturgrammatiken
ausreichend ist. Das Ergebnis dieses Kapitels ist eine Phrasenstrukturgrammatik, die den
\xbar"=Grammatiken (ausgesprochen: „X-bar-Grammatiken``) ähnelt, wie sie in den späten 1970er und frühen
1980er Jahren entwickelt wurden
\citep{Chomsky70a,Jackendoff77a}. Die in diesem Kapitel begründeten Strukturen werden auch in
späteren Kapiteln eine Rolle spielen, aber die lexikalischen Einträge werden wesentlich reichhaltiger
sein: Sie werden Valenzinformationen enthalten, die eine entscheidende Rolle bei der Lizenzierung
syntaktischer Struktur spielen.

Viel Zeit wird auf die Struktur von Nominalphrasen verwendet. Die wesentlichen Einsichten zur Syntax
der \ili{deutsch}en Nominalphrasen lassen sich auf andere \ili{germanisch}e Sprachen übertragen. Spätere
Kapitel werden sich mit den Unterschieden zwischen den \ili{germanisch}en Sprachen bezüglich der Syntax
auf Satzebene befassen.

Dieses Kapitel stützt sich stark auf \citew[Kapitel~2]{MuellerGT-Eng5}, was eine aktualisierte
Übersetzung von \citew[Kapitel~2]{MuellerGTBuch2} ist. Die Kenntnis grundlegender Begriffe wie Wortart
und Konstituententests wird vorausgesetzt. Leser, die das Bedürfnis haben, ihr Wissen in diesen
Bereichen aufzufrischen, seien auf Kapitel~1 dieser Lehrbücher verwiesen.

\section{Symbole und Ersetzungsregeln}
\label{sec-simple-psg}

Wörter lassen sich aufgrund ihrer Flexionseigenschaften und ihrer syntaktischen Distribution einer
bestimmten Wortart zuordnen. So ist \emph{weil} in (\mex{1})
eine Konjunktion\is{Konjunktion}, wohingegen \emph{das} und \emph{dem}
Artikel\is{Artikel} sind und daher zu den Determinatoren\is{Determinator} gezählt werden. Ferner sind \emph{Buch} und \emph{Kind} Nomen\is{Nomen}
und \emph{gibt} ist ein Verb\is{Verb}.
\ea\label{bsp-weil-er-das-buch-dem-mann-gibt}
weil er das Buch dem Kind gibt
\z
Mit den in \citew[Abschnitt~1.3]{MuellerGT-Eng5} eingeführten Konstituententests lässt sich zeigen, dass
einzelne Wörter ebenso wie die Zeichenketten \emph{das Buch} und \emph{dem Kind}
Konstituenten bilden. Diesen werden dann bestimmte Symbole zugewiesen. Da Nomen einen wichtigen Teil
der Phrasen \emph{das Buch} und \emph{dem Kind} bilden, werden diese als \emph{Nominalphrasen} oder
kurz NPs bezeichnet. Das Pronomen \emph{er} kann an denselben Stellen wie volle NPs auftreten und kann
daher auch der Kategorie NP zugeordnet werden.

Die Gruppierung von Konstituenten kann man sich mithilfe von Kästen veranschaulichen und darstellen. Zum Beispiel kann \emph{er das
  Buch dem Kind gibt} wie in Abbildung~\ref{fig-boxes} dargestellt werden.
\begin{figure}
\centering
\TZbox{%
\TZbox{er}
\TZbox{%
       \TZbox{das}
       \TZbox{Buch}}
\TZbox{%
       \TZbox{dem}
       \TZbox{Kind}}
\TZbox{gibt}}
\caption{\label{fig-boxes}Wörter und Phrasen in Kästen}
\end{figure}%

Die oben genannten Kategorien können in dieses Bild integriert werden. Das resultierende Bild ist als
Abbildung~\ref{fig-boxes-pos} dargestellt.
\begin{figure}
\centering
\TZbox{%
\TZbox{\begin{tabular}{@{}l@{}}
       NP\\
       er
       \end{tabular}}
\TZbox{%
       \begin{tabular}{@{}l@{}}
       NP\\
       \TZbox{\begin{tabular}{@{}l@{}}
              Det\\
              das
              \end{tabular}}
       \TZbox{\begin{tabular}{@{}l@{}}
              N\\
              Buch
              \end{tabular}}
       \end{tabular}}
\TZbox{%
       \begin{tabular}{@{}l@{}}
       NP\\
       \TZbox{\begin{tabular}{@{}l@{}}
              Det\\
              dem
              \end{tabular}}
       \TZbox{\begin{tabular}{@{}l@{}}
              N\\
              Kind
              \end{tabular}}
       \end{tabular}}
\TZbox{\begin{tabular}{@{}l@{}}
       V\\
       gibt
       \end{tabular}}}
\caption{\label{fig-boxes-pos}Wörter und Phrasen in Kästen mit Wortartetiketten}
\end{figure}
Kästen mit denselben Etiketten können durch andere Kästen mit demselben Etikett ersetzt werden. Zum Beispiel kann der Kasten
für \emph{dem Kind} durch \emph{dem Mädchen} ersetzt werden. \emph{Er} kann
durch \emph{das Mädchen} ersetzt werden. Das ist sehr intuitiv, aber es ist besser, ein Werkzeug zu haben,
mit dem man tatsächlich Strukturen ableiten kann, die sich als solche Kästen oder als syntaktische
Bäume darstellen lassen, wie man sie aus Einführungskursen kennt. Daher werden wir uns jetzt mit
Phrasenstrukturgrammatiken befassen.

Phrasenstrukturgrammatiken kommen mit Regeln, die festlegen, welche Symbole bestimmten Arten von Wörtern zugewiesen werden und wie diese kombiniert werden, um komplexere
Einheiten zu bilden. Eine einfache Phrasenstrukturgrammatik, mit der man (\mex{0}) analysieren kann, ist in (\mex{1}) angegeben:\footnote{%
	Ich ignoriere die Konjunktion \emph{weil} vorerst. Da die genaue Analyse von
        \ili{deutsch}en Verberst- und Verbzweitsätzen eine Reihe zusätzlicher Annahmen erfordert, beschränken wir uns in diesem Kapitel auf Verbletztsätze.
}$^,$\footnote{\label{fn-np-pron-ps-rule}%
	Die Regel NP $\to$ er mag merkwürdig erscheinen. Wir könnten stattdessen die Regel PersPron $\to$ er annehmen, müssten dann aber eine weitere Regel postulieren,
	die festlegt, dass Personalpronomen volle NPs ersetzen können: NP $\to$ PersPron. Die Regel in (\mex{1}) fasst die beiden genannten Regeln zusammen und besagt,
	dass \emph{er} an Stellen auftreten kann, an denen Nominalphrasen auftreten können.
}
\ea
\label{bsp-grammatik-psg}
\begin{tabular}[t]{@{}l@{ }l}
{NP} & {$\to$ Det N}\\          
{S}  & {$\to$ NP NP NP V}
\end{tabular}\hspace{2cm}%
\begin{tabular}[t]{@{}l@{ }l}
{NP} & {$\to$ er}\\
{Det}  & {$\to$ das}\\
{Det}  & {$\to$ dem}\\
\end{tabular}\hspace{8mm}
\begin{tabular}[t]{@{}l@{ }l}
{N} & {$\to$ Buch}\\
{N} & {$\to$ Kind}\\
{V} & {$\to$ gibt}\\
\end{tabular}
\z
Wir können eine Regel wie NP $\to$\is{$\to$} Det N also so interpretieren, dass eine Nominalphrase, das heißt etwas, dem das Symbol NP zugewiesen ist, aus
einem Determinator (Det) und einem Nomen (N) bestehen kann.

Wir können den Satz in (\mex{-1}) mit der Grammatik in (\mex{0}) folgendermaßen analysieren:
Zunächst nehmen wir das erste Wort des Satzes und prüfen, ob es eine Regel gibt, in der dieses Wort auf der rechten
Seite der Regel vorkommt. Ist dies der Fall, dann ersetzen wir das Wort durch das Symbol auf der linken Seite der Regel. Das geschieht
in den Zeilen 2--4, 6--7 und 9 der Ableitung in (\mex{1}). So wird
zum Beispiel in Zeile~2 \emph{er} durch NP ersetzt.
Wenn zwei oder mehr Symbole gemeinsam auf der rechten Seite einer Regel vorkommen, dann werden alle
diese Wörter durch das Symbol auf der linken Seite ersetzt. Das geschieht in den Zeilen 5, 8 und 10. So
werden zum Beispiel in den Zeilen 5 und 8 Det und N zu NP umgeschrieben.
%\begin{figure}
\ea
%\oneline
{%
\label{bsp-anwendung-grammatik}
\begin{tabular}[t]{@{}r|llllll@{\hspace{1.4cm}}l@{}}
 & \multicolumn{6}{l@{}}{Wörter und Symbole} & angewendete Regeln\\\midrule
 1 & er            & das          & Buch          & dem          & Kind & gibt                \\
 2 & {NP}          & das          & Buch          & dem          & Kind & gibt & {NP $\to$ er}  \\
 3 & NP            & Det          & Buch          & dem          & Kind & gibt & {Det $\to$ das}  \\
 4 & NP            & Det          & N             & dem          & Kind & gibt & {N $\to$ Buch} \\
 5 & NP            &              & NP            & dem          & Kind & gibt & {NP $\to$ Det N}\\
 6 & NP            &              & NP            & Det          & Kind & gibt & {Det $\to$ dem}  \\
 7 & NP            &              & NP            & Det          & N    & gibt & {N $\to$ Kind} \\
 8 & NP            &              & NP            &              & NP   & gibt & {NP $\to$ Det N}\\
 9 & NP            &              & NP            &              & NP   & {V} & {V $\to$ gibt}  \\
10 &               &              &               &              &      & {S} & {S $\to$ NP NP NP V}\\
\end{tabular}}
\z
%\vspace{-\baselineskip}\end{figure}%
In (\mex{0}) sind wir mit einer Wortkette gestartet, und es wurde gezeigt, dass wir die Struktur eines Satzes ableiten können, indem wir die Regeln
einer gegebenen Phrasenstrukturgrammatik anwenden. Wir hätten dieselben Schritte auch in umgekehrter Reihenfolge anwenden können: Ausgehend
vom Satzsymbol S hätten wir die Schritte~9--1 angewendet und wären bei der Wortkette angekommen.
Indem man andere Regeln aus der Grammatik zum Umschreiben der Symbole auswählt, könnte man mit der Grammatik in (\mex{-1})
von S zur Kette \emph{er dem Kind das Buch gibt} gelangen.
Wir können sagen, dass diese Grammatik eine Menge von Sätzen lizenziert (oder generiert)\label{Seite-generiert}.\is{Generative Grammatik}

Die Ableitung in (\mex{0}) kann auch als Baum dargestellt werden. Dies zeigt Abbildung~\vref{er-das-buch-dem-mann-gibt-flat}.
\begin{figure}
\centerline{
\begin{forest}
sm edges
[S
  [NP [er;he] ]
  [NP
    [Det [das;the] ]
    [N [Buch;book] ] 
  ]
  [NP
    [Det [dem;the] ]
    [N [Kind;child] ] 
  ]
  [V [gibt;gives] ]
]
\end{forest}
}
\caption{\label{er-das-buch-dem-mann-gibt-flat}Analyse von \emph{er das Buch dem Kind gibt}}
\end{figure}%
Die Symbole im Baum werden \emph{Knoten}\is{Knoten} genannt. Wir sagen, dass S die NP-Knoten und den V-Knoten unmittelbar dominiert\is{Dominanz}.
Die anderen Knoten im Baum werden ebenfalls von S dominiert, aber nicht unmittelbar dominiert\is{Dominanz!unmittelbare}. Wenn wir über die
Beziehung zwischen Knoten sprechen wollen, ist es üblich, Verwandtschaftsbezeichnungen zu verwenden. In Abbildung~\ref{er-das-buch-dem-mann-gibt-flat} ist S der \emph{Mutterknoten}\is{Knoten!Mutter-}
der drei NP-Knoten und des V-Knotens.
Die NP-Knoten und V sind \emph{Schwestern}\is{Knoten!Schwester-} oder \emph{Töchter}\is{Knoten!Tochter-}, da
sie denselben Mutterknoten haben.\footnote{\emph{Elternknoten}\is{Knoten!Eltern-} und \emph{Kindknoten}\is{Knoten!Kind-} sind alternative
  Begriffe. Ich verwende hier \emph{Mutter} und \emph{Tochter}, da diese Terminologie auch in
  Formalisierungen der später entwickelten Theorie verwendet wird.
}
Wenn ein Knoten zwei Töchter\is{Knoten!Tochter-} hat, dann haben wir eine binär\is{binär} verzweigende Struktur\is{Verzweigung}\is{Verzweigung!binäre}.
Wenn es genau eine Tochter gibt, dann haben wir eine unär\is{unär}\is{Verzweigung!unäre} verzweigende
Struktur. Zwei Wörter oder Phrasen heißen \emph{adjazent}\is{Adjazenz},
wenn sie direkt nebeneinander stehen.

Phrasenstrukturregeln werden in linguistischen Publikationen oft weggelassen. Stattdessen entscheiden sich die Autoren für Baumdiagramme oder die kompakte äquivalente Klammernotation
wie in (\mex{1}).
\ea
{}[\sub{S} [\sub{NP} er] [\sub{NP} [\sub{Det} das] [\sub{N} Buch]]  [\sub{NP} [\sub{Det} dem] [\sub{N} Kind]] [\sub{V} gibt]]
\z
Trotzdem sind es die grammatischen Regeln/""Schemata, die eigentlich wichtig sind, da diese grammatisches Wissen darstellen, das unabhängig von spezifischen Strukturen ist.
Auf diese Weise können wir die Grammatik in (\ref{bsp-grammatik-psg}) verwenden, um den Satz
in (\mex{1}) zu analysieren oder zu generieren, der sich von (\ref{bsp-weil-er-das-buch-dem-mann-gibt}) in der Abfolge der Objekte unterscheidet:
\ea
(weil) er dem Kind das Buch gibt
\z
Die Regeln zum Ersetzen von Determinatoren und Nomen werden einfach in einer anderen Reihenfolge angewendet als in (\ref{bsp-weil-er-das-buch-dem-mann-gibt}). Statt den ersten Det durch \emph{das} und das erste Nomen durch \emph{Buch} zu ersetzen, wird der erste Det durch \emph{dem} und das erste Nomen durch \emph{Kind} ersetzt.

An dieser Stelle sollte ich darauf hinweisen, dass die Grammatik in (\ref{bsp-grammatik-psg}) nicht die einzig mögliche Grammatik für den Beispielsatz in
(\ref{bsp-weil-er-das-buch-dem-mann-gibt}) ist. Es gibt unendlich\label{page-unendlich-viele-grammatiken} viele mögliche Grammatiken, mit denen man
diese Art von Sätzen analysieren könnte (siehe \citealt[Kapitel~2,
Übung~1]{MuellerGT-Eng}). Eine andere mögliche Grammatik ist die folgende:

\ea\label{psg-binaer}
\begin{tabular}[t]{@{}l@{ }l@{}}
NP & $\to$ Det N  \\
V  & $\to$ NP V\\
\end{tabular}\hspace{2cm}%
\begin{tabular}[t]{@{}l@{ }l}
{NP}  & {$\to$ er}\\
{Det} & {$\to$ das}\\
{Det} & {$\to$ dem}\\
\end{tabular}\hspace{8mm}
\begin{tabular}[t]{@{}l@{ }l}
{N} & {$\to$ Buch}\\
{N} & {$\to$ Kind}\\
{V} & {$\to$ gibt}\\
\end{tabular}
\z
Diese Grammatik lizenziert nur binär verzweigende Strukturen, wie in Abbildung~\vref{er-das-buch-dem-mann-gibt-bin} gezeigt.
\begin{figure}
\centerline{
\begin{forest}
sm edges
[V
  [NP [er;he] ]
  [V
    [NP
      [Det [das;the] ]
      [N [Buch;book] ] ]
    [V
      [NP
        [Det [dem;the] ]
        [N [Kind;child] ] ]
      [V [gibt;gives] ] ] ] ]
\end{forest}
}
\caption{\label{er-das-buch-dem-mann-gibt-bin}Analyse von \emph{er das Buch dem Kind gibt} mit einer
  binär verzweigenden Struktur}
\end{figure}%
Da die Regel V $\to$ NP, V rekursiv ist, können beliebig viele NPs mit einem V kombiniert werden. Das Ergebnis
einer NP-V-Kombination ist ein V, das wiederum als Tochter auf der rechten Seite der Regel verwendet werden kann.

Sowohl die Grammatik in (\mex{0}) als auch die in (\ref{bsp-grammatik-psg}) ist zu ungenau.
Wenn wir zusätzliche lexikalische Einträge für \emph{ich} und \emph{den} (Akkusativ) in unsere
Grammatik aufnehmen, dann würde die Grammatik fälschlicherweise
die ungrammatischen Sätze in (\mex{1}b--d) lizenzieren:\footnote{\label{fn-ex-das-kind-erwartet}%
	Mit der Grammatik in (\ref{psg-binaer}) haben wir außerdem das zusätzliche Problem, dass wir nicht bestimmen können, wann eine Äußerung vollständig ist,
	da das Symbol V für alle Kombinationen von V und NP verwendet wird. Daher können wir mit dieser Grammatik auch
        die Sätze in (i) analysieren, sofern wir die entsprechenden Wörter ins Lexikon aufnehmen:

\eal
\ex[*]{
der Delphin erwartet
}
\ex[*]{
des Kindes er das Buch dem Kind gibt
}
\zl
Die Anzahl der von einem Verb geforderten Argumente muss irgendwie in der Grammatik repräsentiert werden. In Kapitel~\ref{chap-valence} werden wir genau sehen,
wie die Selektion von Argumenten durch ein Verb (Valenz) in HPSG erfasst wird.
}
\eal
\ex[]{
er das Buch dem Kind gibt
}
\ex[*]{
ich das Buch dem Kind gibt
}
\ex[*]{
er das Buch den Kind gibt
}
\ex[*]{
er den Buch dem Kind gibt
}
\zl
In (\mex{0}b) wurde die Subjekt-Verb-Kongruenz\is{Kongruenz|(} verletzt. Mit anderen Worten: \emph{ich} und \emph{gibt} passen nicht zusammen.
(\mex{0}c) ist ungrammatisch, weil die Kasusanforderungen des Verbs nicht erfüllt sind: \emph{gibt} verlangt ein Dativobjekt. Schließlich ist (\mex{0}d) ungrammatisch,
weil zwischen Determinator und Nomen keine Kongruenz besteht. Es ist nicht möglich, \emph{den}, das maskulin ist und Akkusativ trägt,
und \emph{Buch} zu kombinieren, weil \emph{Buch} Neutrum ist. Da die Genusmerkmale
dieser beiden Elemente nicht übereinstimmen, können die Elemente nicht kombiniert werden.

Im Folgenden werden wir überlegen, wie wir unsere Grammatik ändern müssten, um sie daran zu hindern, die Sätze in (\mex{0}b--d) zu lizenzieren.
Wenn wir die Subjekt-Verb-Kongruenz erfassen wollen, dann müssen wir die folgenden sechs Fälle im \ili{Deutsch}en abdecken, da das Verb mit dem
Subjekt sowohl in Person\is{Person} (1, 2, 3) als auch in Numerus\is{Numerus} (Sg, Pl) kongruieren muss\is{Kongruenz}:
\eal\jamwidth=7cm\relax%\settowidth\jamwidth{(3, sg)}
\ex
Ich schlafe. \jambox{(1, sg)}
\ex
Du schläfst. \jambox{(2, sg)}
\ex
Er/sie/es schläft. \jambox{(3, sg)}
\ex
Wir schlafen. \jambox{(1, pl)}
\ex
Ihr schlaft. \jambox{(2, pl)}
\ex
Sie schlafen. \jambox{(3, pl)}
\zl
Es ist möglich, diese Beziehungen mit grammatischen Regeln zu erfassen, indem man
die Anzahl der verwendeten Symbole erhöht. Statt der Regel S $\to$ NP NP NP V können wir
die folgende verwenden:
\ea
\begin{tabular}[t]{@{}l@{ }l@{~~}l@{~~}l@{~~}l}
S  & $\to$ NP\_1\_sg & NP & NP & V\_1\_sg\\
S  & $\to$ NP\_2\_sg & NP & NP & V\_2\_sg\\
S  & $\to$ NP\_3\_sg & NP & NP & V\_3\_sg\\
S  & $\to$ NP\_1\_pl & NP & NP & V\_1\_pl\\
S  & $\to$ NP\_2\_pl & NP & NP & V\_2\_pl\\
S  & $\to$ NP\_3\_pl & NP & NP & V\_3\_pl\\
\end{tabular}
\z
Das würde bedeuten, dass wir jeweils sechs verschiedene Symbole für Nominalphrasen und Verben benötigen sowie sechs Regeln statt einer.\is{Kongruenz|)}

Um die Kasus\isi{Kasus}zuweisung durch das Verb zu erfassen, können wir Kasusinformationen auf analoge Weise in die Symbole aufnehmen. Wir würden dann
Regeln wie die folgenden erhalten:
\ea
\label{ditrans-ps-regeln}
\begin{tabular}[t]{@{}l@{ }l@{~~}l@{~~}l@{~~}l}
S  & $\to$ NP\_1\_sg\_nom & NP\_dat & NP\_acc & V\_1\_sg\_nom\_dat\_acc\\
S  & $\to$ NP\_2\_sg\_nom & NP\_dat & NP\_acc & V\_2\_sg\_nom\_dat\_acc\\
S  & $\to$ NP\_3\_sg\_nom & NP\_dat & NP\_acc & V\_3\_sg\_nom\_dat\_acc\\
S  & $\to$ NP\_1\_pl\_nom & NP\_dat & NP\_acc & V\_1\_pl\_nom\_dat\_acc\\
S  & $\to$ NP\_2\_pl\_nom & NP\_dat & NP\_acc & V\_2\_pl\_nom\_dat\_acc\\
S  & $\to$ NP\_3\_pl\_nom & NP\_dat & NP\_acc & V\_3\_pl\_nom\_dat\_acc\\
\end{tabular}
\z
Da es nötig ist, zwischen Nominalphrasen in vier Kasus zu unterscheiden, haben wir insgesamt sechs Symbole für NPs im Nominativ und drei Symbole für NPs mit
anderen Kasus. Da Verben zu den NPs passen müssen, das heißt, da wir zwischen Verben, die drei Argumente selegieren, und solchen, die nur eins oder zwei selegieren, unterscheiden müssen (\mex{1}),
müssen wir die Anzahl der Symbole erhöhen, die wir für Verben annehmen.\is{Valenz}
\eal
\ex[]{
Aicke schläft.
}
\ex[*]{
Aicke schläft das Buch.
}
\ex[]{
Aicke kennt das Buch.
}
\ex[*]{
Aicke kennt.
}
\zl
In den obigen Regeln ist die Information über die Anzahl der von einem Verb geforderten Argumente in
den atomaren Symbolen enthalten, \zb `nom\_dat\_acc'.

Um die Determinator-Nomen-Kongruenz\is{Kongruenz|(} in (\mex{1}) zu erfassen, müssen wir Informationen über Genus\is{Genus} (fem, mas, neu),
Numerus\is{Numerus} (Sg, Pl), Kasus\is{Kasus} (Nom, Gen, Dat, Akk) und die Flexionsklassen (stark, schwach) aufnehmen.\footnote{%
Dies sind Flexionsklassen für Adjektive, die auch für einige Nomen wie \emph{Beamter},
\emph{Verwandter}, \emph{Gesandter} relevant sind.
%For more on adjective classes see page~\pageref{page-Flexionsklasse-Wunderlich}.%
}
\eal\settowidth\jamwidth{(Inflectional class)}
\ex
der Mann, die Frau, das Buch \jambox{(Genus)}
\ex
das Buch, die Bücher \jambox{(Numerus)}
\ex
des Buches, dem Buch \jambox{(Kasus)}
\ex\is{Flexionsklasse}
ein Beamter, der Beamte \jambox{(Flexionsklasse)}
\zl
\largerpage
Statt der Regel NP $\to$ Det N müssen wir Regeln wie die in (\mex{1}) verwenden.\footnote{%
  Um die Dinge einfach zu halten, enthalten diese Regeln keine Information bezüglich der Flexionsklasse.
}
(\mex{1}) zeigt die Regeln für Nominalphrasen im Nominativ. Wir würden analoge Regeln für Genitiv,
Dativ und Akkusativ benötigen. Wir würden dann 24 Symbole für Determinatoren ($3*2*4$), 24 Symbole für Nomen und
24 Regeln statt einer benötigen. Wenn die Flexionsklasse berücksichtigt wird, verdoppelt sich die Anzahl der Symbole und die
Anzahl der Regeln.\is{Kongruenz|)}
\ea
%\resizebox{\linewidth}{!}{
\begin{tabular}[t]{@{}l@{ }l@{~~}l}
NP\_3\_sg\_nom  & $\to$ Det\_fem\_sg\_nom & N\_fem\_sg\_nom \\
NP\_3\_sg\_nom  & $\to$ Det\_mas\_sg\_nom & N\_mas\_sg\_nom \\
NP\_3\_sg\_nom  & $\to$ Det\_neu\_sg\_nom & N\_neu\_sg\_nom \\
NP\_3\_pl\_nom  & $\to$ Det\_fem\_pl\_nom & N\_fem\_pl\_nom \\
NP\_3\_pl\_nom  & $\to$ Det\_mas\_pl\_nom & N\_mas\_pl\_nom \\
NP\_3\_pl\_nom  & $\to$ Det\_neu\_pl\_nom & N\_neu\_pl\_nom \\
\end{tabular}
\z
% moved before the example because of layout
% (\mex{0}) shows the rules for nominative noun phrases. We would need analogous rules for genitive,
% dative, and accusative. We would then require 24 symbols for determiners ($3*2*4$), 24 symbols for nouns and
% 24 rules rather than one. If inflection class is taken into account, the number of symbols and the
% number of rules doubles.\is{agreement|)} 

\section{Erweiterung der PSG um Merkmale}
\label{sec-PSG-Merkmale}

Phrasenstrukturgrammatiken, die nur atomare Symbole verwenden, sind problematisch, da sie bestimmte Generalisierungen nicht erfassen können.
Wir als Linguisten können erkennen, dass NP\_3\_sg\_nom für eine Nominalphrase steht, weil es die Buchstaben NP enthält.
In formaler Hinsicht ist dieses Symbol jedoch genauso wie jedes andere Symbol in der Grammatik, und wir können die Gemeinsamkeiten
aller für NPs verwendeten Symbole nicht erfassen. Ferner erfassen unstrukturierte Symbole nicht die Tatsache, dass die Regeln in (\mex{0})
alle etwas gemeinsam haben. In formaler Hinsicht ist das Einzige, was die Regeln gemeinsam haben, dass es ein Symbol auf der
linken Seite der Regel und zwei auf der rechten gibt.

Wir können dieses Problem lösen, indem wir Merkmale einführen, die den Kategoriesymbolen zugewiesen werden und es daher erlauben, die Werte
solcher Merkmale in unsere Regeln aufzunehmen. Zum Beispiel können wir die Merkmale Person, Numerus und Kasus für das Kategoriesymbol
NP annehmen. Für Determinatoren und Nomen würden wir ein zusätzliches Merkmal für Genus und eines für
die Flexionsklasse annehmen. (\mex{1}) zeigt zwei Regeln, die um die jeweiligen Werte in Klammern ergänzt sind:\footnote{%
  In den folgenden Kapiteln werden Merkmal-Wert-Strukturen verwendet. In diesen Strukturen haben wir immer
Paare aus einem Merkmalnamen und einem Merkmalwert. In einem solchen Rahmen ist die Reihenfolge der Werte nicht
wichtig, da jeder Wert durch den entsprechenden Merkmalnamen eindeutig identifiziert wird. Da wir in Schemata wie (\mex{0}) keinen Merkmalnamen
haben, ist die Reihenfolge der Werte wichtig.
}

\ea
\begin{tabular}[t]{@{}l@{ }l}
NP(3,sg,nom)  & $\to$ Det(fem,sg,nom) N(fem,sg,nom)\\
NP(3,sg,nom)  & $\to$ Det(mas,sg,nom) N(mas,sg,nom)\\
\end{tabular}
\z
\largerpage
Würden wir Variablen statt der Werte in (\mex{0}) verwenden, so erhielten wir Regelschemata wie das
in (\mex{1}):
\ea
\label{Regel-mit-Variablen}
\begin{tabular}[t]{@{}l@{ }l@{ }l}
NP({3},{Num},{Case}) & $\to$ & Det(Gen,{Num},{Case}) N(Gen,{Num},{Case})\\
\end{tabular}
\z
Die Werte der Variablen sind hier nicht wichtig. Wichtig ist, dass sie übereinstimmen. Damit
dies funktioniert, ist es wichtig, dass die Werte geordnet sind; das heißt, in der Kategorie eines Determinators steht das Genus immer an erster Stelle, der Numerus
an zweiter und so weiter. Der Wert des Personmerkmals (die erste Position in der NP(3,Num,Case)) wird durch die Regel auf `3' festgelegt. Diese
Art von Beschränkungen der Werte kann natürlich im Lexikon festgelegt werden:
\ea
\begin{tabular}[t]{@{}l@{ }l}
NP(3,sg,nom)  & $\to$ es\\
Det(mas,sg,gen)  & $\to$ des\\
\end{tabular}
\z

\noindent
Die Regeln in (\ref{ditrans-ps-regeln}) können zu einem einzigen Schema wie in (\mex{1}) zusammengefasst werden:
\ea
\label{ditrans-schema}
\begin{tabular}[t]{@{}l@{ }l@{ }l}
S  & $\to$ & NP({Per1},{Num1},{nom}) \\
   &       & NP(Per2,Num2,{dat})\\
   &       & NP(Per3,Num3,{acc})\\
   &       & V({Per1},{Num1},ditransitive)\\
\end{tabular}
\z
Die Identifikation von Per1 und Num1 am Verb und am Subjekt stellt sicher, dass Subjekt-Verb-Kongruenz vorliegt.
Für die anderen NPs sind die Werte dieser Merkmale irrelevant. Der Kasus dieser NPs wird explizit festgelegt.
\is{Phrasenstrukturgrammatik|)}

% \section{Semantics}
% \label{sec-PSG-Semantik}

% In the introductory chapter and the previous sections, we have been dealing with syntactic aspects
% of language and the focus will remain very much on syntax for the remainder of this book. It is,
% however, important to remember that we use language to communicate, that is, to transfer information
% about certain situations, topics or opinions. If we want to accurately explain our capacity for
% language, then we also have to explain the meanings that our utterances have. To this end, it is
% necessary to understand their syntactic structure, but this alone is not enough. Furthermore,
% theories of language acquisition that only concern themselves with the acquisition of syntactic
% constructions are also inadequate. The syntax"=semantics interface\is{syntax"=semantics interface}
% is therefore important and every grammatical theory has to say something about how syntax and
% semantics interact. In the following, I will show how we can combine phrase structure rules with
% semantic information. To represent meanings, I will use first"=order predicate logic and
% $\lambda$"=calculus\is{$\lambda$"=calculus}. Unfortunately, it is not possible to provide a detailed
% discussion of the basics of logic so that even readers without prior knowledge can follow all the
% details, but the simple examples discussed here should be enough to provide some initial
% insights into how syntax and semantics interact and furthermore, how we can develop a linguistic
% theory to account for this.

% To show how the meaning of a sentence is derived from the meaning of its parts, we will consider (\mex{1}a). We
% assign the meaning in (\mex{1}b) to the sentence in (\mex{1}a). 
% \eal
% \ex\label{Bsp-Max-schlaeft}
% \gll Max schläft.\\
%      Max sleeps\\
% \glt `Max is sleeping.'
% \ex\label{Bsp-schlafen-max} 
% \relation{schlafen}(\relation{max})
% \zl
% Here, we are assuming \relation{schlafen} to be the meaning of  \emph{schläft} `sleeps'. We use prime symbols to indicate
% that we are dealing with word meanings and not actual words. At first glance, it may not seem that we have really gained anything
% by using \relation{schlafen} to represent the meaning of (\mex{0}a), since it is just another form of the verb \emph{schläft} `sleeps'.
% It is, however, important to concentrate on a single verb form as inflection is irrelevant when it comes to meaning. We can see this by comparing the 
% examples in (\mex{1}a) and (\mex{1}b):

% \eal
% \ex 
% \gll Jeder Junge schläft.\\
%      every boy sleeps\\
% \glt `Every boy sleeps.'
% \ex 
% \gll Alle Jungen schlafen.\\
%      all boys sleep\\
% \glt `All boys sleep.'	 
% \zl

% \noindent
% To enhance readability I use \ili{English} translations of the predicates in semantic representations
% from now on.\footnote{%
%   Note that I do not claim that \ili{English} is suited as representation language for semantic relations
%   and concepts that can be expressed in other languages.
% }
% So the meaning of (\mex{-1}a) is represented as (\mex{1}) rather then (\mex{-1}b):

% \ea
% \label{sleep-max}
% \relation{sleep}(\relation{max})
% \z
% When looking at the meaning in (\mex{-0}), we can consider which part of the meaning comes from each word.
% It seems relatively intuitive that \relation{max} comes from \emph{Max}, but the trickier question is what exactly
% \emph{schläft} `sleeps' contributes in terms of meaning. If we think about what characterizes a `sleeping' event, we
% know that there is typically an individual who is sleeping. This information is part of the meaning of the verb \emph{schlafen}
% `to sleep'. The verb meaning does not contain information about the sleeping individual, however, as this verb can be used
% with various subjects:
%  \eal
% \ex 
% \gll Paul schläft.\\
%      Paul sleeps\\
% \glt `Paul is sleeping.'
% \ex 
% \gll Mio schläft.\\
%      Mio sleeps\\
% \glt `Mio is sleeping.'
% \ex 
% \gll Xaver schläft.\\
%      Xaver sleeps\\
% \glt `Xaver is sleeping.'
% \zl
% We can therefore abstract away from any specific use of \relation{sleep} and instead of, for example, \relation{max} in (\ref{sleep-max}), we
% use a variable (\zb $x$). This $x$ can then be replaced by \relation{paul}, \relation{mio} or \relation{xaver} in a given sentence. To allow us
% to access these variables in a given meaning, we can write them with a $\lambda$ in front. Accordingly, \emph{schläft} `sleeps' will have
% the following meaning:
% \ea
% $\lambda x~\relation{sleep}(x)$
% \z
% %
% The step from (\ref{sleep-max}) to (\mex{0}) is referred to as \emph{lambda abstraction}\is{lambda"=abstraction@$\lambda$"=abstraction}.
% The combination of the expression (\mex{0}) with the meaning of its arguments happens in the following way: we remove the $\lambda$ and the 
% corresponding variable and then replace all instances of the variable with the meaning of the
% argument. If we combine (\mex{0}) and \relation{max} as in (\mex{1}),
% we arrive at the meaning in (\ref{sleep-max}), namely \relation{sleep}(\relation{max}). 
% \ea
% $\lambda x~\relation{sleep}(x)$ \relation{max}
% \z
% The process is called $\beta$"=reduction\is{beta"=reduction@$\beta$"=reduction} or 
% $\lambda$"=conversion\is{lambda"=conversion@$\lambda$"=conversion}. To show this further, let us consider an example with a transitive verb. The sentence
% in (\mex{1}a) has the meaning given in (\mex{1}b):
% \eal
% \ex\label{Bsp-Max-mag-Lotte} 
% \gll Max mag Lotte.\\
%      Max likes Lotte\\
% \glt `Max likes Lotte.'
% \ex \relation{like}(\relation{max}, \relation{lotte})
% \zl
% The $\lambda$"=abstraction of \emph{mag} `likes' is shown in (\mex{1}):
% \ea
% $\lambda y \lambda x~\relation{like}(x, y)$
% \z
% Note that it is always the first $\lambda$ that has to be used first. The variable $y$ corresponds
% to the object of \emph{mögen} `to like'. For languages like \ili{English} it is assumed that the object forms a verb
% phrase (VP) together with the verb and this VP is combined with the subject. \ili{German} differs from
% \ili{English} in allowing more freedom in constituent order. The problems that result for form meaning
% mappings are solved in different ways by different theories. The respective solutions will be
% addressed in the following chapters.


% If we combine the representation in (\mex{0}) with that of the object \emph{Lotte}, we arrive at (\mex{1}a), and following
% $\beta$"=reduction, (\mex{1}b):
% \eal
% \label{lambda-moegen}
% \ex $\lambda y \lambda x~\relation{like}(x, y) \relation{lotte}$
% \ex $\lambda x~\relation{like}(x, \relation{lotte})$
% \zl

% \noindent
% This meaning can in turn be combined with the subject and we then get (\mex{1}a) and (\mex{1}b) after $\beta$"=reduction:
% \eal
% \ex $\lambda x~\relation{like}(x, \relation{lotte}) \relation{max}$
% \ex \relation{like}(\relation{max}, \relation{lotte})
% \zl

% \begin{sloppypar}
% \noindent
% After introducing lambda calculus, integrating the composition of meaning into our phrase structure rules is simple. A rule for the
% combination of a verb with its subject has to be expanded to include positions for the semantic contribution of the verb, the semantic
% contribution of the subject and then the meaning of the combination of these two (the entire sentence). The complete meaning is the
% combination of the individual meanings in the correct order. We can therefore take the simple rule in (\mex{1}a) and turn it into 
% (\mex{1}b):
% \end{sloppypar}
% \eal
% \ex S $\to$ NP(nom) V
% \ex S(V$'$ NP$'$) $\to$ NP(nom, NP$'$) V(V$'$)
% \zl
% V$'$ stands for the meaning of V and NP$'$ for the meaning of the NP(nom). V$'$ NP$'$ stands for the combination of V$'$ and NP$'$. When analyzing
% (\ref{Bsp-Max-schlaeft}), the meaning of V$'$ is $\lambda x~\relation{sleep}(x)$ and the meaning of NP$'$ is \relation{max}. The combination of V$'$ NP$'$
% corresponds to (\mex{1}a) or after $\beta$"=reduction to (\ref{Bsp-schlafen-max}) -- repeated here as (\mex{1}b):
% \eal
% \ex $\lambda x~\relation{sleep}(x) \relation{max}$
% \ex \relation{sleep}(\relation{max})
% \zl

% \noindent
% For the example with a transitive verb in (\ref{Bsp-Max-mag-Lotte}), the rule in (\mex{1}) can be proposed:
% \ea
% S(V$'$ NP2$'$ NP1$'$) $\to$ NP(nom, NP1$'$) V(V$'$) NP(acc, NP2$'$)
% \z
% The meaning of the verb (V$'$) is first combined with the meaning of the object (NP2$'$) and then with the meaning of the subject (NP1$'$). 

% At this point, we can see that there are several distinct semantic rules for the phrase structure rules above. The hypothesis that we should analyze language
% in this way is called the \emph{rule"=to"=rule hypothesis}\is{rule"=to"=rule hypothesis}
% \citep[\page 184]{Bach76a}. A more general process for deriving the
% meaning of linguistic expression will be presented in Section~\ref{Sec-GPSG-Sem}.

\section{Phrasenstrukturregeln für einige Aspekte der deutschen Syntax}

Während die Bestimmung der direkten Konstituenten eines Satzes relativ einfach ist, da wir uns wegen der
einigermaßen flexiblen Abfolge der Konstituenten im \ili{Deutsch}en stark auf den Verschiebetest stützen können, ist es schwieriger, die Teile der Nominalphrase zu identifizieren. Das ist das Problem,
auf das wir uns in diesem Abschnitt konzentrieren werden. Um die in Abschnitt~\ref{sec-xbar} zu besprechenden Annahmen über die \xbar"=Syntax zu motivieren,
werden wir auch Präpositionalphrasen besprechen.


\subsection{Nominalphrasen}
\label{sec-psg-np}

\largerpage
Bislang haben wir eine relativ einfache Struktur für Nominalphrasen angenommen: Unsere Regeln besagen, dass eine Nominalphrase aus einem Determinator und einem
Nomen besteht. Nominalphrasen können eine deutlich komplexere Struktur als (\mex{1}a) haben. Das zeigen die folgenden Beispiele in (\mex{1}):
\eal
\label{Beispiele-NP-Adjunkte}
\ex
ein Buch
\ex
\label{ex-ein-Buch-das-wir-kennen}
ein Buch, das wir kennen
\ex
\label{ex-ein-Buch-aus-Japan}
ein Buch aus Japan
\ex
ein interessantes Buch
\ex
ein Buch aus Japan, das wir kennen
\ex
ein interessantes Buch aus Japan
\ex
ein interessantes Buch, das wir kennen
\ex
ein interessantes Buch aus Japan, das wir kennen
\zl

\noindent
Neben Determinatoren und Nomen können Nominalphrasen auch Adjektive, Präpositionalphrasen und Relativsätze enthalten.
Die zusätzlichen Elemente in (\mex{0}) sind Adjunkte\is{Adjunkt|(}. Sie schränken die Menge der Objekte ein, auf die sich die Nominalphrase
bezieht. Während sich (\mex{0}a) auf eine Entität bezieht, die die Eigenschaft hat, ein Buch zu sein, muss der Referent von (\mex{0}b)
zusätzlich die Eigenschaft haben, uns bekannt zu sein.

Unsere bisherigen Regeln für Nominalphrasen kombinierten einfach ein Nomen und einen Determinator und können daher nur zur
Analyse von (\mex{0}a) verwendet werden. Die Frage, vor der wir nun stehen, ist, wie wir diese Regel modifizieren oder welche zusätzlichen Regeln wir
annehmen müssten, um die anderen Nominalphrasen in (\mex{0}) zu analysieren. Zusätzlich zu Regel (\mex{1}a) könnte man
eine Regel wie die in (\mex{1}b) vorschlagen.\footnote{%
	Siehe \citew[\page 238]{Eisenberg2004a} für die Annahme flacher Strukturen in Nominalphrasen.
}$^,$\footnote{%
	Es gibt natürlich weitere Merkmale wie Genus und Numerus, die Teil aller in diesem Abschnitt
	besprochenen Regeln sein sollten. Ich habe diese im Folgenden der einfacheren Darstellung halber weggelassen.
}
%\todostefan{These footnotes have to be blocked from moving to the next page.}
% reformulating one line to be shorter plus enlarging the page by two lines did the trick.
\eal
\ex NP $\to$ Det N
\ex NP $\to$ Det A N
\zl
\largerpage[2]
Diese Regel würde es uns jedoch immer noch nicht erlauben, Nominalphrasen wie (\mex{1}) zu analysieren:
\ea
\label{Beispiel-alle-weitern-schlagkraeftigen-Argumente}
alle weiteren schlagkräftigen Argumente
\z
Um (\mex{0}) analysieren zu können, benötigen wir eine Regel wie (\mex{1}):
\ea 
NP $\to$ Det A A N
\z
Es ist immer möglich, die Anzahl der Adjektive in einer Nominalphrase zu erhöhen, und eine Obergrenze für
Adjektive festzulegen wäre völlig willkürlich. Selbst wenn wir uns für die folgende Abkürzung entscheiden, gibt es weiterhin Probleme:

\ea
NP $\to$ Det A* N
\z
Der Stern\is{*} in (\mex{0}) steht für eine beliebige Anzahl von Iterationen. Daher umfasst (\mex{0}) Regeln ohne Adjektive
ebenso wie solche mit einem, zwei oder mehr.

Das Problem ist, dass laut der Regel in (\mex{0}) Adjektive und Nomen keine Konstituente bilden und wir daher nicht erklären können, warum Koordination
in (\mex{1}) trotzdem möglich ist:
\ea
\label{ex-alle-grossen-Seeelefanten-und}
alle [[großen Seeelefanten] und [grauen Eichhörnchen]]
\z
Wenn wir annehmen, dass Koordination die Kombination von zwei oder mehr Wortketten mit denselben syntaktischen Eigenschaften umfasst, dann müssten wir annehmen,
dass Adjektiv und Nomen eine Einheit bilden.

%The noun phrases with adjectives discussed thus far can be captured by the following rules:
%
% Shorter for layout:
%The following rules capture the noun phrases with adjectives discussed thus far:
Die Regeln in (\mex{1}) erfassen die bislang besprochenen Nominalphrasen mit Adjektiven:
\eal
\label{NP-Regeln}
\ex NP $\to$ Det \nbar
\ex\label{NP-Regeln-Adj} \nbar $\to$ A \nbar
\ex\label{NP-Regeln-Nbar-N} \nbar $\to$ N
\zl

\largerpage
\noindent
Diese Regeln besagen Folgendes: Eine Nominalphrase besteht aus einem Determinator und einem nominalen Element (\nbar). Dieses nominale Element
kann aus einem Adjektiv und einem nominalen Element (\mex{0}b) bestehen oder nur aus einem Nomen (\mex{0}c). Da \nbar auch auf der rechten Seite
der Regel in (\mex{0}b) steht, können wir diese Regel mehrfach anwenden und so Nominalphrasen mit mehreren Adjektiven wie
(\ref{Beispiel-alle-weitern-schlagkraeftigen-Argumente}) erfassen. Abbildung~\vref{Abbildung-Adjektive-in-NP} zeigt die Struktur einer Nominalphrase
ohne Adjektiv und die einer Nominalphrase mit einem oder zwei Adjektiven.
\begin{figure}
%\hfill%
\scalebox{.9}{%
\begin{forest}
sm edges
[NP
   [Det [ein;a] ]
   [\nbar
      [N [Eichhörnchen;squirrel] ] ] ]
\end{forest}}
\hfill
\scalebox{.9}{%
\begin{forest}
sm edges
[NP
   [Det [ein;a] ]
   [\nbar
      [A [graues;gray] ]
      [\nbar
        [N [Eichhörnchen;squirrel] ] ] ] ]
\end{forest}}
%
\hfill
\scalebox{.9}{%
\begin{forest}
sm edges
[NP
  [Det [ein;a] ]
    [\nbar
    [A [großes;big] ]
       [\nbar
       [A [graues;gray] ]
         [\nbar
         [N [Eichhörnchen;squirrel] ] ] ] ] ]
\end{forest}}
%\hfill
%\mbox{}
%
\caption{\label{Abbildung-Adjektive-in-NP}Nominalphrasen mit unterschiedlicher Anzahl von Adjektiven}
\end{figure}%
Das Adjektiv \emph{grau} schränkt die Menge der Referenten der Nominalphrase ein. Wenn wir ein
zusätzliches Adjektiv wie \emph{groß} annehmen, dann bezieht sie sich nur auf diejenigen Eichhörnchen, die grau
und auch groß sind. Diese Art von Nominalphrasen kann in Kontexten wie dem folgenden verwendet werden:

\ea
\label{Beispiel-Iteration-Adjektive}
A: Alle grauen Eichhörnchen sind groß.\\
B: Nein, ich habe ein kleines graues Eichhörnchen gesehen.
\z
Wir beobachten, dass dieser Diskurs mit \emph{Aber alle kleinen grauen Eichhörnchen
  sind krank} und einer entsprechenden Antwort fortgesetzt werden kann. Die Möglichkeit,
  noch mehr Adjektive in Nominalphrasen wie \emph{ein kleines graues Eichhörnchen} zu haben,
  wird in unserem Regelsystem in (\mex{-1}) erfasst. In der Regel
  (\ref{NP-Regeln-Adj}) kommt \nbar sowohl auf der linken als auch auf der rechten
  Seite der Regel vor. Diese Art von Regel wird als \emph{rekursiv}\is{Rekursion} bezeichnet.
\is{Adjunkt|)}

\largerpage
Wir haben nun eine raffinierte kleine Grammatik entwickelt, mit der man Nominalphrasen mit
adjektivischen Modifikatoren analysieren kann. Dadurch wird der Kombination aus einem Adjektiv und einem Nomen Konstituenten\-status
verliehen. Man könnte sich an dieser Stelle fragen, ob es nicht sinnvoll wäre, auch anzunehmen, dass Determinatoren und
Adjektive eine Konstituente bilden, da wir auch die folgende Art von Nominalphrasen haben:
\ea
diese schlauen und diese neugierigen Eichhörnchen
\z
%\largerpage
Hier haben wir es jedoch mit einer anderen Struktur zu tun. Zwei volle NPs wurden
koordiniert\is{Koordination}, und ein Teil des ersten Konjunkts wurde nicht ausgesprochen.\footnote{
  Man beachte, dass man nicht behaupten kann, dass das zweite Konjunkt in (\ref{ex-alle-grossen-Seeelefanten-und})
  eine volle NP ist und \emph{alle} nur nicht ausgesprochen wird. Wenn der Determinator in \ili{deutsch}en
  NPs weggelassen wird, brauchen wir eine andere Flexion:
\eal
\ex[]{
Alle grauen Eichhörnchen sind groß.
}
\ex[]{
Graue Eichhörnchen sind groß.
}
\ex[*]{
Grauen Eichhörnchen sind groß.
}
\zl
Was in (\ref{ex-alle-grossen-Seeelefanten-und}) koordiniert wird, sind also zwei Adjektiv-Nomen-Kombinationen,
und das Ergebnis wird mit einem Determinator kombiniert.
}

\ea
diese schlauen \st{Eichhörnchen} und diese neugierigen Eichhörnchen
\z
Ähnliche Phänomene findet man auf Satzebene (\mex{1}a) und sogar auf Wortebene (\mex{1}b):
\eal
\ex
dass Conny dem Kind das Buch \st{gibt} und Aicke der Frau die Schallplatte gibt
\ex
be- und ent"=laden
\zl
Koordination ist ein komplexes Phänomen. Siehe \citew{AC2021a} für einen Überblick.
% Dass in (\mex{-1}) wirklich keine normale symmetrische Koordination vorliegt, sieht man, wenn man
% (\mex{-1}) mit (\mex{1}) vergleicht:
% \ea
% diese schlauen Frauen und klugen Männer
% \z
% Mit (\mex{0}) verweist man auf eine Gruppe, die aus schlauen Frauen und klugen Männern besteht,
% wohingegen man mit (\mex{-2}) auf zwei Gruppen verweist, nämlich


Bislang haben wir besprochen, wie wir Adjektive idealerweise in unsere Regeln für die Struktur von Nominalphrasen integrieren können.
Andere Adjunkte wie Präpositionalphrasen (\ref{ex-ein-Buch-aus-Japan}) oder Relativsätze (\ref{ex-ein-Buch-das-wir-kennen}) können analog zu Adjektiven mit \nbar kombiniert werden:
\eal
\ex\label{xbar-PP-Adjunkt-an-N} \nbar $\to$ \nbar PP
\ex \nbar $\to$ \nbar Relativsatz
\zl
Mit diesen Regeln und denen in (\ref{NP-Regeln}) ist es möglich -- unter Annahme der entsprechenden Regeln für PPs und
Relativsätze -- alle Beispiele in (\ref{Beispiele-NP-Adjunkte}) zu analysieren.

(\ref{NP-Regeln}c) besagt, dass es möglich ist, dass \nbar aus einem einzigen Nomen besteht. Eine weitere wichtige Regel wurde noch nicht
besprochen: Wir benötigen eine weitere Regel, um Nomen wie \emph{Vater}, \emph{Sohn} oder \emph{Bild},
so"=genannte \emph{relationale Nomen}\is{Nomen!relationales}, mit ihren Argumenten zu kombinieren. Beispiele dafür finden sich in (\mex{1}a--b).
(\mex{1}c) ist ein Beispiel für eine Nominalisierung eines Verbs mit seinem Argument:
\eal
\label{Beispiele-NP-relationale-Nomina}
\ex
der Vater von Peter
\ex
das Bild vom Gleimtunnel
\ex
das Kommen der Installateurin
\zl
\noindent
Die Regel, die wir zur Analyse von (\ref{Beispiele-NP-relationale-Nomina}a,b) benötigen, ist in
(\mex{1}) angegeben:
%\todostefan{Martin: It is not said why relational nouns should behave in a special way.}
\ea
\nbar $\to$ N PP
\z
%
Abbildung~\ref{Abbildung-NP-mit-PP-Argument} zeigt zwei Strukturen mit PP"=Argumenten. Der Baum auf der rechten Seite enthält außerdem ein zusätzliches PP"=Adjunkt, das durch
die Regel in (\ref{xbar-PP-Adjunkt-an-N}) lizenziert wird.
\begin{figure}
\centerfit{%
\begin{forest}
sm edges
[NP
 [Det [das;the] ]
 [\nbar
   [N [Bild;picture] ]
   [PP [vom Gleimtunnel;of.the Gleimtunnel,roof ] ] ] ]
\end{forest}%
\hspace{2em}%
\begin{forest}
sm edges
[NP
  [Det [das;the] ]
  [\nbar
    [\nbar
      [N [Bild;picture] ]
      [PP [vom Gleimtunnel;of.the Gleimtunnel,roof ] ] ] 
    [PP [im Gropiusbau;in.the Gropiusbau,roof ] ] ] ]
\end{forest}}
\caption{\label{Abbildung-NP-mit-PP-Argument}Kombination eines Nomens mit dem PP"=Komplement
  \emph{vom Gleimtunnel} rechts mit einem Adjunkt-PP}
\end{figure}%

\largerpage
Zusätzlich zu den zuvor besprochenen NP-Strukturen gibt es weitere Strukturen, in denen der
Determinator oder das Nomen fehlt.
Nomen können durch Ellipse weggelassen werden. (\mex{1}) gibt ein Beispiel für Nominalphrasen, in denen ein Nomen, das kein Komplement verlangt,
weggelassen wurde. Die Beispiele in (\mex{2}) zeigen NPs, in denen nur ein Determinator und ein Komplement des Nomens realisiert wurde,
aber nicht das Nomen selbst. Der Unterstrich markiert die Position, an der das Nomen normalerweise auftreten würde.
\eal
\label{ex-nounless-np}
\ex
ein interessantes \_
\ex
ein neues interessantes \_

\ex
ein interessantes \_ aus  Japan
\ex
ein interessantes \_, das  wir kennen
\zl

\eal
\label{ex-nounless-np-relational-noun}
\ex
(Nein, nicht der Vater von Klaus), der \_ von Peter war gemeint.
%\itdopt{that statt the one?}
\ex
(Nein, nicht das Bild von der Stadtautobahn), das \_ vom Gleimtunnel war beeindruckend.
\ex
(Nein, nicht das Kommen des Tischlers), das \_ der Installateurin ist wichtig.
\zl
Im \ili{Englisch}en muss oft das Pronomen
\emph{one} an der entsprechenden Stelle verwendet werden,\footnote{%
  Siehe \citet[Section~4.12]{FLGR2012a} für \ili{englisch}e Beispiele ohne das
  Pronomen \emph{one}.
} aber im \ili{Deutsch}en wird das Nomen\is{Nomen}
einfach weggelassen.
%\todostefan{ich habe hier mit absicht ``often'' hinzugefügt, weil in den obengenannten beispielen
%man nicht in allen Fällen ``one'' benutzen kann, vgl. (\mex{0}c)}
In Phrasenstrukturgrammatiken kann dies durch eine so"=genannte \emph{Epsilon-Produktion}\is{Epsilon-Produktion}\is{leeres Element} beschrieben werden.
Diese Regeln ersetzen ein Symbol durch nichts (\mex{1}a). Die Regel in (\mex{1}b) ist eine äquivalente Variante, die für den Begriff \emph{Epsilon-Produktion} verantwortlich ist:
\eal
\label{np-epsilon}
\ex N $\to$
\ex N $\to$ $\epsilon$
\zl 

\noindent
Die entsprechenden Bäume sind in Abbildung~\vref{Abbildung-NP-ohne-Nomen} dargestellt.
\begin{figure}
\hfill
\begin{forest}
sm edges
[NP
  [Det [ein;an] ]
  [\nbar
    [A [interessantes;interesting] ]
    [\nbar
      [N [\trace ] ] ] ] ]
\end{forest}
\hfill
\begin{forest}
sm edges
[NP
  [Det [das;the] ]
  [\nbar
    [N [\trace] ]
    [PP [vom Gleimtunnel;of.the Gleimtunnel, roof] ] ] ]
\end{forest}
\hfill%
\mbox{}
\caption{\label{Abbildung-NP-ohne-Nomen}Nominalphrasen ohne overten Kopf}
\end{figure}%
Kehren wir zu Kästen wie dem in Abbildung~\ref{fig-boxes-pos} zurück, so entsprechen die Regeln in (\mex{0}) leeren Kästen mit denselben Etiketten wie die Kästen
gewöhnlicher Nomen. Wie wir zuvor überlegt haben, ist der tatsächliche Inhalt der Kästen unwichtig, wenn
man die Frage betrachtet, wo wir sie einfügen können. Zum Beispiel können die Nominalphrasen in (\ref{Beispiele-NP-Adjunkte})
in denselben Sätzen auftreten. Ebenso verhält sich der leere Nomenkasten wie einer mit einem echten Nomen: Wenn wir
den leeren Kasten nicht öffnen, werden wir den Unterschied zu einem gefüllten Kasten nicht bemerken können.

Es ist nicht nur möglich, das Nomen aus Nominalphrasen wegzulassen, sondern auch der Determinator kann in bestimmten Kontexten unrealisiert bleiben.
(\mex{1}) zeigt Nominalphrasen im Plural\is{Plural}:
\eal
\ex
Bücher
\ex
Bücher, die wir kennen
\ex
interessante Bücher
\ex
interessante Bücher, die  wir kennen
\zl
Der Determinator kann auch im Singular weggelassen werden, wenn das Nomen ein Massennomen\is{Nomen!Massen-} bezeichnet:

\eal
\ex
Getreide
\ex
Getreide, das gerade gemahlen wurde
\ex
frisches Getreide
\ex
frisches Getreide, das gerade gemahlen wurde
\zl
Schließlich können sowohl der Determinator als auch das Nomen weggelassen werden:
\eal
\ex
Ich lese interessante.
\ex
Dort drüben steht frisches, das gerade gemahlen wurde.
\zl
Abbildung~\vref{Abbildung-NP-ohne-Det} zeigt die entsprechenden Bäume.

\begin{figure}
\hfill
\begin{forest}
sm edges
[NP
  [Det [\trace] ]
  [\nbar
    [N [Bücher;books] ] ] ]
\end{forest}
\hfill
\begin{forest}
sm edges
[NP
  [Det [\trace] ]
  [\nbar
    [A [interessante;interesting] ]
    [\nbar
      [N [\trace] ] ] ] ]
\end{forest}
\hfill
\mbox{}
\caption{\label{Abbildung-NP-ohne-Det}Nominalphrasen ohne overten Determinator}
\end{figure}%

\largerpage
Es ist nötig, zwei weitere Anmerkungen zu den bislang entwickelten Regeln zu machen: Bisher habe ich
immer von Adjektiven gesprochen. Es ist jedoch möglich, sehr komplexe Adjektivphrasen in prä"=nominaler Position zu haben.
Dabei kann es sich um Adjektive mit Komplementen (\mex{1}a,b) oder um adjektivische Partizipien (\mex{1}c,d) handeln:

\eal
\ex
der seiner Frau treue Mann
\ex
der auf seinen Sohn stolze Mann
\ex
der seine Frau liebende Mann
\ex
der von seiner Frau geliebte Mann
\zl
Vor diesem Hintergrund muss die Regel (\ref{NP-Regeln-Adj}) folgendermaßen modifiziert werden:
\ea
\label{NP-Regeln-AP}
\nbar $\to$ AP \nbar
\z
Eine Adjektivphrase (AP) kann aus einer NP und einem Adjektiv, einer PP und einem Adjektiv oder nur aus einem Adjektiv bestehen:
\eal
\ex AP $\to$ NP A
\ex AP $\to$ PP A
\ex AP $\to$ A
\zl
Es gibt zwei Unzulänglichkeiten, die sich aus den bislang entwickelten Regeln ergeben. Das sind die Regeln für Adjektive
oder Nomen ohne Komplemente in (\mex{0}c) sowie (\ref{NP-Regeln-Nbar-N}) -- hier wiederholt als (\mex{1}):
\ea
\nbar $\to$ N
\z
Wenn wir diese Regeln anwenden, dann erzeugen wir unär verzweigende Teilbäume, das heißt Bäume mit einer Mutter, die
nur eine Tochter hat. (Siehe Abbildung~\ref{Abbildung-NP-ohne-Det} für ein Beispiel dafür.) Wenn wir die
Parallele zu den Kästen aufrechterhalten, würde dies bedeuten, dass es einen Kasten gibt, der einen weiteren Kasten enthält, nämlich denjenigen mit
dem relevanten Inhalt.

Im Prinzip hindert uns nichts daran, diese Information direkt in den größeren Kasten zu legen. Statt
der Regeln in (\mex{1}) verwenden wir einfach die Regeln in (\mex{2}):
\eal
\ex A $\to$ kluge
\ex N $\to$ Mann
\zl
\eal
\label{Lexikon-Projektion}
\ex AP $\to$ kluge
\ex \nbar $\to$ Mann
\zl
(\mex{0}a) besagt, dass \emph{kluge} dieselben Eigenschaften wie eine volle Adjektivphrase hat, insbesondere, dass es nicht
mit einem Komplement kombiniert werden kann. Dies ist parallel zur Kategorisierung des Pronomens \emph{er} als NP in den Grammatiken
(\ref{bsp-grammatik-psg}) und (\ref{psg-binaer}).

\largerpage[2]
Die Zuweisung der Kategorie \nbar an Nomen, die kein Komplement verlangen, hat den Vorteil, dass wir nicht erklären müssen, warum die Analyse in (\mex{1}b) ebenso möglich ist wie
(\mex{1}a), obwohl es keinen Bedeutungsunterschied gibt.
\eal
\ex
{}[\sub{NP} einige [\sub{\nbar} kluge [\sub{\nbar} [\sub{\nbar} [\sub{N} Frauen ] und  [\sub{\nbar} [\sub{N} Männer ]]]]]]
\ex
{}[\sub{NP} einige [\sub{\nbar} kluge [\sub{\nbar} [\sub{N} [\sub{N} Frauen ] und [\sub{N} Männer
]]]]]
\zl
%
In (\mex{0}a) haben zwei Nomen zu \nbar projiziert und wurden dann durch Koordination verbunden. Das Ergebnis der Koordination
zweier Konstituenten derselben Kategorie ist immer eine neue Konstituente mit dieser Kategorie. Im Fall von (\mex{0}a) ist dies
ebenfalls \nbar. Diese Konstituente wird dann mit dem Adjektiv und dem Determinator kombiniert. In (\mex{0}b) wurden die Nomen selbst
koordiniert. Das Ergebnis davon ist immer eine weitere Konstituente, die dieselbe Kategorie wie ihre Teile hat. In diesem Fall
wäre dies N. Dieses N wird zu \nbar und wird dann mit dem Adjektiv kombiniert. Wären Nomen, die keine Komplemente verlangen,
als \nbar statt als N kategorisiert, so hätten wir nicht das Problem der unechten
Ambiguitäten\is{Ambiguität!unechte}.\footnote{
Äußerungen in natürlicher Sprache sind oft mehrdeutig. Zum Beispiel hat der folgende Satz zwei
Lesarten.
\ea
Unbekannte haben Mittwochabend bei einer FDP-Wahlkampfveranstaltung mit FDP-Chef Guido Westerwelle Farbbeutel geworfen. (taz, 21.5.2004, p.\,7)
\z
Die beiden Lesarten entsprechen zwei verschiedenen Strukturen. In der ersten Lesart hängt die \emph{mit}-PP
an die Wahlkampfveranstaltung an, was bedeutet, dass Guido Westerwelle auf der Veranstaltung war. In der zweiten
Lesart modifiziert die PP das Verb \emph{geworfen}, was einer Bedeutung entspricht, in der Guido
Westerwelle gemeinsam mit Unbekannten Farbbeutel geworfen hat. Das ist normale Ambiguität. Was Linguisten
üblicherweise vermeiden wollen, ist unechte Ambiguität: Fälle, in denen wir dieselbe Semantik, aber zwei
verschiedene syntaktische Strukturen haben.
}
Die Struktur in (\mex{1}) zeigt die einzig mögliche Analyse.
\ea
{}[\sub{NP} einige [\sub{\nbar} kluge [\sub{\nbar} [\sub{\nbar} Frauen ] und [\sub{\nbar} Männer
]]]]
\z

\subsection{Präpositionalphrasen}
\label{Abschnitt-PP-Syntax}

Verglichen mit der Syntax von Nominalphrasen ist die Syntax von Präpositional\is{Präposition|(}phrasen (PPs) relativ geradlinig. PPs bestehen normalerweise
aus einer Präposition und einer Nominalphrase, deren Kasus durch diese Präposition bestimmt wird. Wir können dies mit der folgenden
Regel erfassen:
\ea
\label{Regel-PP-einfach}
PP $\to$ P NP
\z
Diese Regel muss natürlich auch Informationen über den Kasus der NP enthalten. Ich habe dies der einfacheren Darstellung halber weggelassen, so wie ich es
bei den NP"=Regeln und AP"=Regeln oben getan habe.

Die Duden-Grammatik \citep[§~1300]{Duden2005-Authors} bietet Beispiele wie die in (\mex{1}), die zeigen, dass bestimmte Präpositionalphrasen
dazu dienen, den semantischen Beitrag der Präposition näher zu bestimmen, indem sie zum Beispiel ein Maß angeben:
\eal
\ex\label{Beispiel-Schritt-vor-dem-Abgrund}
{}[[Einen Schritt] vor dem Abgrund] blieb er stehen.
\ex
{}[[Kurz] nach dem Start] fiel die Klimaanlage aus.
\ex
{}[[Schräg] hinter der Scheune] ist ein Weiher.
\ex
{}[[Mitten] im Urwald] stießen die Forscher auf einen alten Tempel.
\zl
Um die Sätze in (\mex{0}a,b) zu analysieren, könnte man die folgenden Regeln in (\mex{1}) vorschlagen:

\eal
\ex PP $\to$ NP PP
\ex PP $\to$ AP PP
\zl
Diese Regeln kombinieren eine PP mit einer Maßangabe. Die resultierende Konstituente ist wiederum eine PP. Es ist
möglich, mit diesen Regeln die Präpositionalphrasen in (\mex{-1}a,b) zu analysieren, aber sie erlauben es uns leider auch,
diejenigen in (\mex{1}) zu analysieren:
\eal
\ex[*]{
[\sub{PP} einen Schritt [\sub{PP} kurz [\sub{PP} vor dem Abgrund]]]
}
\ex[*]{
[\sub{PP} kurz [\sub{PP} einen Schritt [\sub{PP} vor dem Abgrund]]]
}
\zl
Beide Regeln in (\mex{-1}) wurden zur Analyse der Beispiele in (\mex{0}) verwendet. Da das Symbol PP sowohl auf der linken
als auch auf der rechten Seite der Regeln vorkommt, können wir die Regeln in beliebiger Reihenfolge und beliebig oft anwenden.
% Fn: Semantik hilft nicht.

Wir können diesen unerwünschten Nebeneffekt vermeiden, indem wir die zuvor angenommenen Regeln umformulieren:
\eal
\ex PP $\to$ NP \pbar
\ex PP $\to$ AP \pbar
\ex PP $\to$ \pbar\label{Regel-PP-P}
\ex \pbar $\to$ P NP
\zl
Regel (\ref{Regel-PP-einfach}) wird zu (\mex{0}d). Die Regel in (\mex{0}c) besagt, dass eine PP aus \pbar bestehen kann.
Abbildung~\vref{Abbildung-PP} zeigt die Analyse von (\mex{1}) mit (\mex{0}c) und (\mex{0}d) sowie die Analyse
eines Beispiels mit einem Adjektiv an der ersten Position nach den Regeln in (\mex{0}b) und (\mex{0}d):
\ea
vor dem Abgrund
\z
\begin{figure}
\hfill
\begin{forest}
sm edges
[PP
  [\pbar
    [P [vor;before] ]
    [NP [dem Abgrund;the abyss, roof] ] ] ]
\end{forest}
\hfill
\begin{forest}
sm edges
[PP
  [AP [kurz;shortly,roof] ]
  [\pbar
    [P [vor;before] ]
    [NP [dem Abgrund;the abyss,roof] ] ] ]
\end{forest}
%% \begin{tikzpicture}
%% \tikzset{level 1+/.style={level distance=2\baselineskip}}
%% \tikzset{frontier/.style={distance from root=8\baselineskip}}
%% \Tree[.PP
%%        [.{\pbar}
%%          [.P vor;before ]
%%          [.NP \edge[roof]; {dem Abgrund;the abyss} ] ] ] 
%% \end{tikzpicture}
%% \hfill
%% \begin{tikzpicture}
%% \tikzset{level 1+/.style={level distance=2\baselineskip}}
%% \tikzset{frontier/.style={distance from root=8\baselineskip}}
%% \Tree[.PP
%%        [.AP \edge[roof]; {kurz;shortly} ]
%%        [.{\pbar}
%%          [.P vor;before ]
%%          [.NP \edge[roof]; {dem Abgrund;the abyss} ] ] ] 
%% \end{tikzpicture}
\hfill
\mbox{}
\caption{\label{Abbildung-PP}Präpositionalphrasen mit und ohne Maßangabe}
\end{figure}%

An dieser Stelle fragt sich der aufmerksame Leser vermutlich, warum es in der linken Abbildung von
Abbildung~\ref{Abbildung-PP} keine leere Maßphrase gibt, die man in Analogie zum leeren Determinator in Abbildung~\ref{Abbildung-NP-ohne-Det} erwarten könnte.
Der Grund für den leeren Determinator in Abbildung~\ref{Abbildung-NP-ohne-Det} ist, dass die gesamte Nominalphrase
ohne den Determinator eine ähnliche Bedeutung hat wie diejenigen mit einem Determinator. Die Bedeutung, die normalerweise
vom sichtbaren Determinator beigetragen wird, muss irgendwie in die Struktur der Nominalphrase aufgenommen werden. Dies
kann durch den leeren Determinator mit einem entsprechend spezifizierten Bedeutungsbeitrag geschehen.
% If we
% did not place this meaning in the empty determiner, this would lead to more complicated assumptions about semantic
% combination: we only really require the mechanisms presented in Section~\ref{sec-PSG-Semantik} and these are
% very general in nature. The meaning is contributed by the words themselves and not by any rules. If we were
% to assume a unary branching rule such as that in the left tree in Figure~\ref{Abbildung-PP} instead of the
% empty determiner, then this unary branching rule would have to provide the semantics of the determiner. This
% kind of analysis has also been proposed by some researchers.\todostefan{Martin: obscure} See Chapter~\ref{Abschnitt-Diskussion-leere-Elemente} for
% more on empty elements.

Anders als determinatorlosen NPs fehlt Präpositionalphrasen ohne Grad- oder Maßangabe
keine Bedeutungskomponente für die Komposition. Es ist daher nicht nötig, eine leere Maßangabe anzunehmen, die
irgendwie zur Bedeutung der gesamten PP beiträgt. Daher besagt die Regel in (\ref{Regel-PP-P}), dass eine
Präpositionalphrase aus \pbar besteht, das heißt aus einer Kombination von P und NP.\is{Präposition|)}

\section{\xbart}
%\section{X̅-Theory}
\label{sec-xbar}
\rohead{\thesection~~$\skoverline{X}$ Theory}

Wenn\is{X theory@\xbar theory|(} wir uns die im vorigen Abschnitt formulierten Regeln noch einmal ansehen, sehen wir, dass Köpfe immer
mit ihren Komplementen zu einer neuen Konstituente kombiniert werden (\mex{1}a,b), die dann mit weiteren Konstituenten kombiniert werden kann (\mex{1}c,d):

\eal
\ex \nbar $\to$ N PP
\ex \pbar $\to$ P NP
\ex\label{Regel-NP-Xbar}
    NP $\to$ Det \nbar
\ex PP $\to$ NP \pbar
\zl
%
Grammatiker, die sich mit dem \ili{Englisch}en\il{Englisch} beschäftigten, bemerkten, dass parallele Strukturen für Phrasen verwendet werden können, die Adjektive oder Verben als Kopf haben.
Ich bespreche an dieser Stelle Adjektivphrasen und verschiebe die Diskussion der Verbalphrasen auf
Kapitel~\ref{chap-valence}, da die Annahmen bezüglich der Struktur von Sätzen sowohl im \ili{Deutsch}en
als auch im \ili{Englisch}en von der \xbart abweichen, wie sie heute üblicherweise angenommen wird. Wie im \ili{Deutsch}en können bestimmte Adjektive
im \ili{Englisch}en Komplemente nehmen, mit der wichtigen Einschränkung, dass Adjektivphrasen mit Komplementen diese im \ili{Englisch}en nicht prä"=nominal realisieren können.
(\mex{1}) gibt einige Beispiele für Adjektivphrasen:
\eal
\ex Kim and Sandy are proud.
\ex Kim and Sandy are very proud.
\ex Kim and Sandy are proud of their child.
\ex Kim and Sandy are very proud of their child.
\zl
Anders als bei Präpositionalphrasen sind die Komplemente von Adjektiven normalerweise optional. \emph{proud} kann mit oder ohne PP verwendet werden.
Der Gradausdruck \emph{very} ist ebenfalls optional.

Die Regeln, die wir für diese Analyse benötigen, sind in (\mex{1}) angegeben, mit den entsprechenden Strukturen in Abbildung~\vref{Abbildung-AP}.

\begin{samepage}
\eal
\ex AP $\to$ \abar
\ex AP $\to$ AdvP \abar
\ex \abar $\to$ A PP
\ex \abar $\to$ A
\zl
\end{samepage}

\begin{figure}
\hfill
\begin{forest}
sm edges
[AP
  [\abar
    [A [proud] ] ] ]
\end{forest}
\hfill
\begin{forest}
sm edges
[AP
  [AdvP [very] ]
  [\abar
    [A [proud] ] ] ]
\end{forest}
\hfill
\begin{forest}
sm edges
[AP
  [\abar
    [A [proud] ]
    [PP [of their child,roof] ] ] ]
\end{forest}
\hfill
\begin{forest}
sm edges
[AP
  [AdvP [very] ]
  [\abar
    [A [proud] ]
    [PP [of their child,roof] ] ] ]
\end{forest}
\hfill\mbox{}
\caption{\label{Abbildung-AP}Englische Adjektivphrasen}
\end{figure}%

\noindent
Wie in Abschnitt~\ref{sec-PSG-Merkmale} gezeigt wurde, ist es möglich, über sehr spezifische
Phrasenstrukturregeln zu generalisieren und so zu allgemeineren Regeln zu gelangen. Auf diese Weise werden Eigenschaften wie
Person, Numerus und Genus nicht mehr in den Kategoriesymbolen kodiert, sondern es werden nur einfache
Symbole wie NP, Det und N verwendet. Es ist nur dann nötig, etwas über die Werte
eines Merkmals festzulegen, wenn dies im Kontext einer gegebenen Regel relevant ist. Wir können diese Abstraktion einen Schritt
weiterführen: Statt explizite Kategoriesymbole wie N, V, P und A für lexikalische Kategorien und
NP, VP, PP und AP für phrasale Kategorien zu verwenden, kann man einfach eine Variable für die jeweilige Wortart verwenden und von X und XP sprechen.

Diese Form der Abstraktion findet sich in der so"=genannten \xbart (oder X"=bar-Theorie; der Begriff \emph{bar}
bezieht sich auf den Strich über dem Symbol), die von \citet{Chomsky70a} entwickelt und von
\citet{Jackendoff77a} verfeinert wurde. Diese Form abstrakter Regeln spielt in vielen verschiedenen
Theorien eine wichtige Rolle. Zum Beispiel: Government \& Binding\indexgb \citep{Chomsky81a}, Generalized Phrase
Structure Grammar\indexgpsg \citep{GKPS85a,Uszkoreit87a} und Lexical Functional Grammar\indexlfg
\citep{Bresnan82a-ed,BATW2016a}. In der HPSG\indexhpsg, der in diesem Buch angenommenen Theorie, spielt die \xbar-Theorie ebenfalls
eine Rolle, aber es wurden nicht alle Beschränkungen des \xbar-Schemas übernommen.

\largerpage
(\mex{1}) zeigt eine mögliche Instanziierung der \xbar-Regeln, wobei die Kategorie X anstelle von N verwendet wurde, sowie Beispiele für Wortketten,
die mit diesen Regeln abgeleitet werden können:
%\todostefan{martin meint ich solle hier NP statt Strichnotation verwenden}
\eanoraggedright
\label{psg-xbar-schema}
\oneline{%
\begin{tabular}[t]{@{}l@{\hspace{5mm}}l@{\hspace{5mm}}l@{}}
\xbar"=Regel & \mbox{mit spezifischen Kategorien} & \mbox{Beispielketten}\\[2mm]
$\overline{\overline{\mbox{X}}} \rightarrow \overline{\overline{\mbox{specifier}}}$~~\xbar &
$\overline{\overline{\mbox{N}}} \rightarrow \overline{\overline{\mbox{DET}}}$~~\nbar & \mbox{the [picture of Paris]} \\
$\xbar \rightarrow$ \xbar~~$\overline{\overline{\mbox{adjunct}}}$            & \nbar $\rightarrow$ \nbar~~$\overline{\overline{\mbox{REL\_CLAUSE}}}$ & \mbox{[picture of Paris]}\\
                            &                                              & \mbox{[that everybody knows]}\\
\xbar $\rightarrow \overline{\overline{\mbox{adjunct}}}$~~\xbar            & \nbar $\rightarrow \overline{\overline{\mbox{A}}}$~~\nbar & \mbox{beautiful [picture of Paris]}\\
\xbar $\rightarrow$ \mbox{X}~~$\overline{\overline{\mbox{complement}}}*$   & \nbar $\rightarrow$ \mbox{N}~~$\overline{\overline{\mbox{P}}}$ & \mbox{picture [of Paris]}\\
\end{tabular}}
\z

\noindent
Jede Wortart kann X ersetzen (\zb V, A oder P). Das X ohne Strich steht in den obigen Regeln für einen
lexikalischen Eintrag. Wenn man die Strichebene explizit machen will, dann kann man \xnull schreiben.
Genau wie bei der Regel in (\ref{Regel-mit-Variablen}), wo wir den Kasuswert des
Determinators oder des Nomens nicht spezifiziert, sondern lediglich verlangt haben, dass die Werte auf der rechten Seite der
Regel übereinstimmen, verlangen die Regeln in (\mex{0}), dass die Wortart eines Elements auf der rechten Seite
der Regel (X oder \xbar) mit der des Elements auf der linken Seite der Regel (\xbar oder
$\overline{\overline{\mbox{X}}}$) übereinstimmt.

%\largerpage
Ein lexikalisches Element kann mit allen seinen Komplementen\is{Komplement} kombiniert werden. Der `*'\is{*} in der letzten Regel steht für
eine unbegrenzte Anzahl von Wiederholungen des Symbols, dem er folgt. Ein Spezialfall ist das nullfache Auftreten von Komplementen. In \emph{das Bild} ist kein
PP-Komplement von \emph{Bild} vorhanden, und somit wird N zu \nbar. Das Ergebnis der
Kombination eines lexikalischen Elements mit seinen Komplementen ist eine neue Projektionsebene von X: die Projektionsebene 1, die durch
einen Strich markiert wird. \xbar kann dann mit Adjunkten kombiniert werden. Diese können links oder rechts von \xbar auftreten. Das Ergebnis dieser Kombination ist
weiterhin \xbar, das heißt, die Projektionsebene wird durch die Kombination mit einem Adjunkt nicht verändert\is{Adjunkt}.
Maximale Projektionen werden durch zwei
Striche markiert. Man kann auch XP\is{XP} für eine Projektion von X mit zwei Strichen schreiben.
Eine XP besteht aus einem Spezifikator\is{Spezifikator} und \xbar. Je nach
den theoretischen Annahmen sind von einem Verb abhängige Subjekte Spezifikatoren einer Verbalphrase
(\citealt*[100--103]{SWB2003a}; \citealt[Abschnitt~3.1]{MOe2011a})
% Die argumentieren gegen die IP, aber Berman hat VP -> DP, VP und Haider will auch keinen Spec.
%(\citealp[]{Haider97a};
% Haider95b-u
%\citealp[Section~3.2.2]{Berman2003a}) 
und Determinatoren sind Spezifikatoren in NPs \citep[\page
210]{Chomsky70a}. Ferner werden auch Gradmodifikatoren \citep[\page
210]{Chomsky70a} in Adjektivphrasen und Maßangaben in Präpositionalphrasen zu den Spezifikatoren gezählt.

\largerpage
Nicht"=Kopf-Positionen können nur maximale Projektionen aufnehmen, und daher haben Komplemente, Adjunkte und Spezifikatoren immer zwei Striche. Wie bereits oben erwähnt,
hält sich die HPSG nicht an die \xbart. Zum Beispiel können Argumente Wörter oder intermediäre, nicht vollständige
Phrasen sein (siehe Kapitel~\ref{chap-verbal-complex} über den Verbalkomplex in \ili{germanisch}en SOV-Sprachen).
Abbildung~\vref{Abb-GB-Min-Max} gibt einen Überblick über die minimale und maximale Struktur von Phrasen.
\begin{figure}
\hfill
\begin{forest}
sm edges
[XP
  [\xbar [X] ] ]
\end{forest}
\hfill
\begin{forest}
%where n children=0{}{},
%sm edges
%for tree={parent anchor=south, child anchor=north,align=center,base=bottom}
[XP
  [specifier]
  [\xbar
    [adjunct]
    [\xbar
      [complement] [X] ] ] ]
\end{forest}
\hfill\mbox{}
\caption{\label{Abb-GB-Min-Max}Minimale und maximale Struktur von Phrasen}
\end{figure}%

Manche Kategorien haben keinen Spezifikator oder haben die Option, einen zu haben. Adjunkte sind optional, und daher
müssen nicht alle Strukturen ein \xbar mit einer Adjunkt-Tochter enthalten. Zusätzlich zu der in der rechten
Abbildung gezeigten Verzweigung sind manchmal Adjunkte an XP und Kopf"=Adjunkte\is{Adjunkt!Kopf-}
möglich. Es gibt in (\ref{psg-xbar-schema}) nur eine einzige Regel für Fälle, in denen ein Kopf den
Komplementen vorausgeht, jedoch ist eine Abfolge, in der das Komplement dem Kopf vorausgeht, natürlich ebenfalls
möglich. Das zeigt Abbildung~\ref{Abb-GB-Min-Max}.

Abbildung~\vref{Abb-das-schoene-Bild-von-Paris} zeigt die Analyse der NP-Strukturen \emph{das Bild}
und \emph{das schöne Bild von Paris}. Die NP-Strukturen in Abbildung~\ref{Abb-das-schoene-Bild-von-Paris}
und der Baum für \emph{proud} in Abbildung~\ref{Abbildung-AP} zeigen Beispiele für minimal besetzte Strukturen.
Der linke Baum in Abbildung~\ref{Abb-das-schoene-Bild-von-Paris} ist außerdem ein Beispiel für eine Struktur ohne Adjunkt. Die rechte Struktur
in Abbildung~\ref{Abb-das-schoene-Bild-von-Paris} ist ein Beispiel für die maximal besetzte Struktur:
Spezifikator, Adjunkt und Komplement sind vorhanden.


\begin{figure}%[p]
\hfill
\begin{forest}
sm edges
[NP
  [DetP
    [\detbar
      [Det [das;the] ] ] ]
  [\nbar
    [N [Bild;picture] ] ] ]
\end{forest}
\hfill
\begin{forest}
sm edges
[NP
  [DetP
    [\detbar
      [Det [das;the] ] ] ]
  [\nbar
    [AP
      [\abar
        [A [schöne;beautiful] ] ] ]
    [\nbar
      [N [Bild;picture] ]
      [PP 
        [\pbar
          [P [von;of] ]
          [NP
            [\nbar
              [N [Paris;Paris] ] ] ] ] ] ] ] ]
\end{forest}
%
\hfill\mbox{}
\caption{\label{Abb-das-schoene-Bild-von-Paris}\xbar"=Analyse von \emph{das Bild}
  und \emph{das schöne Bild von Paris}}
\end{figure}%

Die in Abbildung~\ref{Abb-das-schoene-Bild-von-Paris} gegebene Analyse nimmt an, dass alle Nicht"=Köpfe in einer Regel
Phrasen sind. Man muss daher annehmen, dass es eine Determinatorphrase gibt, selbst wenn der Determinator nicht mit anderen Elementen kombiniert wird.
Die unäre Verzweigung von Determinatoren ist nicht elegant, aber sie ist konsistent.\footnote{%
	Für eine alternative Version der \xbar-Theorie, die keine ausgearbeitete Struktur für Determinatoren annimmt, siehe \citew{Muysken82a}.
}
Die unären Verzweigungen für die NP \emph{Paris} in Abbildung~\ref{Abb-das-schoene-Bild-von-Paris} mögen ebenfalls etwas merkwürdig erscheinen, aber sie werden tatsächlich
plausibler, wenn man komplexere Nominalphrasen betrachtet:
\eal
\ex
das Paris der dreißiger Jahre
\ex
die Maria aus Hamburg
\zl
Unäre Projektionen sind etwas unelegant, aber das sollte uns hier nicht allzu sehr kümmern, da wir
in der Diskussion der lexikalischen Einträge in (\ref{Lexikon-Projektion}) bereits
gesehen haben, dass unär verzweigende Knoten größtenteils vermieden werden können und dass es in der Tat wünschenswert ist,
solche Strukturen zu vermeiden. Andernfalls erhält man unechte Ambiguitäten\is{Ambiguität!unechte}. In den folgenden
Kapiteln zeige ich, wie die HPSG (light) Determinatoren sowie Nominal-, Adjektiv- und Verbalphrasen analysieren kann, ohne
unäre Regeln anzunehmen. Statt also die Strukturen in Abbildung~\ref{Abb-das-schoene-Bild-von-Paris} anzunehmen,
werden die viel einfacheren in Abbildung~\ref{Abb-das-schoene-Bild-von-Paris-HPSG} verwendet.
\begin{figure}%[p]
\hfill
\begin{forest}
sm edges
[NP
  [Det [das;the] ]
  [\nbar
    [Bild;picture] ] ]
\end{forest}
\hfill
\begin{forest}
sm edges
[NP
  [Det [das;the] ]
  [\nbar
    [AP
      [schöne;beautiful] ]
    [\nbar
      [N [Bild;picture] ]
      [PP 
        [P [von;of] ]
        [NP
          [Paris;Paris] ] ] ] ] ]
\end{forest}
%
\hfill\mbox{}
\caption{\label{Abb-das-schoene-Bild-von-Paris-HPSG}HPSG-Analyse von \emph{das Bild}
  und \emph{das schöne Bild von Paris}}
\end{figure}%


% Furthermore, other \xbar~theoretical assumptions will not be shared by several theories discussed in this book. In particular, the assumption that non"=heads always have
% to be maximal projections\is{projection!maximal} will be disregarded. \citet{Pullum85a} and
% \citet{KP90a} have shown that the respective theories are not necessarily less restrictive
% than theories which adopt a strict version of the \xbar theory. See also the discussion in \citew[Section~13.1.2]{MuellerGT-Eng5}.
\is{X theory@\xbar theory|)}
%\ohead{\headmark}

%\clearpage  
%\bigskip
\questions{

\begin{enumerate}
\item Warum sind Phrasenstrukturgrammatiken, die nur atomare Kategorien verwenden, für die Beschreibung natürlicher Sprachen unzureichend?
\item Geben Sie unter Annahme der Grammatik in (\ref{psg-binaer}) an, welche Schritte (das Ersetzen von Symbolen) man durchführen muss, um zum Symbol
	 V im Satz (\mex{1}) zu gelangen.
\ea
er das Buch dem Kind gibt
\z
Ihre Antwort sollte der Analyse in (\ref{bsp-anwendung-grammatik}) ähneln.
%\pagebreak

% \item Give a representation of the meaning of (\mex{1}) using predicate logic:
% \eal
% \ex 
% \gll Ulrike kennt Hans.\\
%      Ulrike knows Hans\\
% \ex 
% \gll Joshi freut sich.\\
%      Joshi is.happy \REFL{}\\
% \glt `Joshi is happy.'
% \zl
\end{enumerate}}


\exercises{
Dieses Kapitel ist großteils identisch mit Kapitel~2 von \citet{MuellerGT-Eng5}. Da der Fokus
dieses Buches sich von dem des Lehrbuchs über grammatische Theorie unterscheidet, habe ich beschlossen, hier einen anderen
Satz von Übungen bereitzustellen. Wer Interesse daran hat, mehr Übungen zu bearbeiten, kann zusätzlich das
Lehrbuch zur grammatischen Theorie konsultieren. Es ist bei Language Science Press erschienen und daher Open Access, das
heißt, es ist frei verfügbar.
\begin{enumerate}
\item\label{exercise-NP-PSG}
Zeichnen Sie Bäume für die folgenden Phrasen. Sie dürfen das Symbol NP für Eigennamen und \nbar für
  Nomen, die keine Komplemente verlangen, verwenden (wie in Abbildung~\ref{Abb-das-schoene-Bild-von-Paris-HPSG}).
\eal
\ex
eine Stunde vor der Ankunft des Zuges
\ex
kurz nach der Ankunft in Paris
\ex
das ein Lied singende Kind aus dem Allgäu
\zl
\item Schreiben Sie eine Phrasenstrukturgrammatik für die Beispielsätze in (\mex{0}). Verwenden Sie Kategorien wie
  Det, A, N und P für die Phrasenstrukturregeln.

\item Verwenden Sie die Online-Version von SWI-Prolog\footnote{%
  \url{https://swish.swi-prolog.org/}, 2020-06-07.
}
um Ihre Grammatik mit dem Computer zu testen. Einzelheiten zur Notation finden Sie im \ili{englisch}en Wikipedia-Eintrag für Definite Clause Grammar
(DCG)\is{Definite Clause Grammar (DCG)}.\footnote{%
\url{https://en.wikipedia.org/wiki/Definite_clause_grammar}, 2020-06-07.
}
\end{enumerate}}


\furtherreading{
Die Erweiterung von Phrasenstrukturgrammatiken um Merkmale wurde bereits 1963 von \citet{Harman63a} vorgeschlagen.

Die in diesem Kapitel besprochene Phrasenstrukturgrammatik für Nominalphrasen deckt einen großen Teil der Syntax
von Nominalphrasen ab, kann aber bestimmte NP-Strukturen nicht erklären. Ferner hat sie das Problem, das Übung~3 von \citet[Kapitel~2]{MuellerGT-Eng5}
aufzeigen soll. Eine Diskussion dieser Phänomene und eine Lösung im Rahmen der HPSG findet sich
in \citew{Netter98a}. Für eine Diskussion der Frage, ob Det oder N der Kopf in nominalen
Strukturen ist, siehe \citet{MuellerHeadless} und \citet{MyPM2021a}. \citet{VanEynde2021a} ist ein Überblick über die Arbeiten zur NP in der HPSG.

Die Diskussion der Integration semantischer Information in Phrasenstrukturgrammatiken war sehr kurz. Eine detaillierte Diskussion der Prädikatenlogik und
ihrer Integration in Phrasenstrukturgrammatiken -- sowie eine Diskussion des Quantorenskopus\is{Skopus} -- findet sich in \citew{BB2005a}.}


% reset to show the chapter heading
\lehead{\headmark}
\ohead{\headmark}
%      <!-- Local IspellDict: de_DE -->
